% sec-05: the SBIR route against the five-year program, and the delta.
% Table 5.  No figure: the arithmetic is a table and drawing it twice would add
% nothing a reader of Table 5 does not already have.
\section{The SBIR Route Against the Five-Year Program}

The five letters of \S2 ask two different sizes of question. Letters 1 and 2 ask
about a route that can begin within a receipt cycle. Letters 3 and 4 ask about a
route sized to the whole program. This section states the arithmetic that
separates them, and it derives nothing: every figure below is reused from the
capitalization plan \cite{capitalplan2026} and the ten-application work
\cite{tenapps2026}.

\begin{apptable}
\begin{tabularx}{\textwidth}{>{\raggedright\arraybackslash}p{3.8cm} >{\raggedright\arraybackslash}p{2.4cm} >{\raggedright\arraybackslash}X}
\headrow \thead{Quantity} & \thead{Value} & \thead{What it is, and where it is set}\\
\midrule
Program, five years, direct & \$3,500,000 & The five-year ask in the ten-application work, not re-derived here \\
Program, per year, direct & \$700,000 & The same frame, annualized \\
SBIR Phase I, total cost & \$306,000 & Nine months, five milestones \\
SBIR Phase II, total cost & \$1,300,000 & Twenty-four months, seven milestones \\
SBIR route, total cost & \$1,606,000 & Thirty-three months, the two phases summed \\
Direct work inside that award & \$1,396,000 & The route's total less the fee and the indirect allowance \\
Delta the route does not buy & \$2,104,000 & The program's direct cost less the direct work the award funds \\
Private capital behind the firewall & \$5,900,000 & The private position in the capitalization plan \\
Private to federal leverage & 3.67 to 1 & Against the federal annex's 3 to 1 target \\
\end{tabularx}
\end{apptable}
\vspace{-0.60cm}
\tabcap{Table 5. The nine figures every claim in this packet reconciles to, with\\
the source of each stated, and with nothing on this page re-derived.}

\subsection{What the delta means for the letters}

The \$2,104,000 line is why letters 1 and 2 are not the whole day. An SBIR route
at \$1,606,000 over thirty-three months is real, and is the only federal
mechanism in the set written for a firm rather than for a person or an
institution \cite{sbapolicydirective}. It is also insufficient on its own for the
five-year program, and saying so in the letter that asks for it is cheaper than
being found out later.

The private position exists to close that gap, and it sits behind a conflict
firewall rather than beside the award. The chief executive is not a clinical
investigator on the trial this program describes: he is chief executive of the
sponsor and holds the whole of it, so every trigger in 21 CFR 54.2 would fire if
he held both roles, and the investigator role therefore sits wholly with the site
\cite{ecfr2026part54,fdafindisclosure2013}.

\subsection{The eligibility constraint the private side is written against}

Under 13 CFR 121.702, majority ownership by multiple venture capital operating
companies, hedge funds, or private equity firms changes SBIR eligibility and at
some agencies requires a separate authorization \cite{ecfr2026sbirownership}. A
convertible instrument that has not converted does not change the answer today; a
priced round that transfers majority ownership does.

That is why the SBA Company Registry answer in
\appfile{funding/auto-fund/02Sep26/forms} is re-checked at every financing event
rather than annually, and why no instrument is signed before its effect on that
answer is written down. Day 2 sets out three candidate instruments side by side
for exactly that reason.

\subsection{The policy frame the federal letters are written against}

The July 2026 report to the President asks the federal research portfolio to
prioritize the individual scientist over legacy institutions and names the SBIR
clause in three separate chapters as the one mechanism written for a firm
\cite{kratsios2026sciencegoldenage}, while its annexed budget memorandum asks
agencies for long-duration grants and non-federal cost share
\cite{kratsiosvought2026fy2028priorities}. Two executive orders name robotics and
physical artificial intelligence as instruments of discovery and set an
evidentiary standard for federally supported science
\cite{eo14363genesismission,eo14303goldstandard}. Every letter in \S2 is written
inside that frame and none argues with it.
